\chapter{Kết luận và hướng phát triển}
\label{chap:conclusion}

\section{Tổng kết đề tài}

Đồ án đã hoàn thành việc nghiên cứu, thiết kế và hiện thực hệ thống \textbf{Socrates Code Tutor} --- một hệ sinh thái gia sư lập trình AI tích hợp phương pháp Socratic, bao gồm:

\begin{enumerate}
    \item \textbf{VS Code Extension} hoạt động ở hai chế độ Free Mode và Classroom Mode, cho phép sinh viên chat với AI tutor, xem bài tập, nộp bài và nhận phản hồi trực tiếp trong IDE.
    \item \textbf{Backend API} dựa trên FastAPI với kiến trúc LangGraph AI Runtime, điều phối 5 workflow AI chuyên biệt: StudentMentorGraph, TeacherInsightGraph, StudentCoachGraph, ExerciseDraftGraph, và ReverseTeachingGraph.
    \item \textbf{Web Portal} cho giảng viên (quản lý lớp, bài tập, dashboard phân tích) và sinh viên (lịch sử học tập, chatbot điều hướng).
    \item \textbf{Hệ thống chấm bài tự động} dựa trên Judge0 Sandbox với khả năng cô lập an toàn.
    \item \textbf{Tinh chỉnh 3 mô hình ngôn ngữ lớn} (Gemma-4, Qwen2.5-Coder, Qwen2.5-Instruct) bằng kỹ thuật LoRA trên benchmark DebugEval và Socratic Debugging.
\end{enumerate}

\section{Các kết quả đã đạt được}

\subsection{Kết quả nghiên cứu}

\begin{itemize}
    \item Qwen Coder + LoRA đạt cải thiện đồng đều trên cả 4 task DebugEval: Bug Localization +27pp, Bug Identification +18pp, Code Repair +10pp.
    \item Gemma-4 + LoRA đạt cải thiện Code Repair vượt trội 4.6 lần (từ 8.94\% lên 41.30\%).
    \item Sau fine-tuning, cả hai mô hình hội tụ về hiệu năng tương đồng, gợi ý tầm quan trọng của chất lượng dữ liệu.
    \item Mô hình Socratic Tutor đạt Question Rate 96.1\%, Answer Concealment 100\%, On-Topic Relevance 4.93/5, Helpfulness 4.49/5.
\end{itemize}

\subsection{Kết quả kỹ thuật}

\begin{itemize}
    \item Hoàn thành hệ thống end-to-end từ Extension đến Backend, Web, Judge Server, AI Runtime.
    \item Triển khai thành công trên AWS EC2 với Docker Compose.
    \item Các mô hình LoRA adapter được công bố trên Hugging Face.
\end{itemize}

\section{Đóng góp của đề tài}

\begin{enumerate}
    \item \textbf{Đóng góp khoa học:} Cung cấp phân tích so sánh định lượng hiệu quả LoRA fine-tuning trên Code Debugging và Socratic Tutoring, bao gồm phân tích specialization trade-off và convergence behavior.
    \item \textbf{Đóng góp kỹ thuật:} Xây dựng hoàn chỉnh một hệ thống giáo dục AI tích hợp IDE, có kiến trúc rõ ràng và khả năng mở rộng.
    \item \textbf{Đóng góp sư phạm:} Chứng minh tính khả thi của việc ứng dụng phương pháp Socratic vào AI tutor với cơ chế Anti-Solution Guard, đạt tỷ lệ che giấu đáp án 100\%.
    \item \textbf{Tài nguyên mở:} Các mô hình LoRA adapter và pipeline đánh giá được công bố phục vụ nghiên cứu.
\end{enumerate}

\section{Hướng phát triển}

\subsection{Mở rộng mô hình AI}

\begin{itemize}
    \item Nghiên cứu multi-task fine-tuning để khắc phục hiện tượng specialization trade-off ở Gemma-4.
    \item Mở rộng bộ dữ liệu Socratic với các cuộc hội thoại bằng tiếng Việt.
    \item Thử nghiệm QLoRA (4-bit quantization) để giảm yêu cầu VRAM.
    \item Tích hợp Retrieval-Augmented Generation (RAG) để AI tham khảo tài liệu giáo trình khi hướng dẫn.
\end{itemize}

\subsection{Mở rộng tính năng hệ thống}

\begin{itemize}
    \item Hỗ trợ thêm các IDE khác (JetBrains, Zed).
    \item Tích hợp bài tập từ nhiều nguồn (Codeforces, LeetCode, bài tập tùy chỉnh).
    \item Gamification: Hệ thống điểm, huy hiệu, bảng xếp hạng để tăng động lực học tập.
    \item Voice-based tutoring: Hỗ trợ giọng nói cho sinh viên khuyết tật thị giác.
\end{itemize}

\subsection{Đánh giá mở rộng}

\begin{itemize}
    \item Thực hiện user study với quy mô lớn (1 lớp $\sim$ 50 sinh viên) trong 1 học kỳ.
    \item So sánh hiệu quả học tập giữa nhóm sử dụng Socratic AI và nhóm sử dụng ChatGPT thông thường.
    \item Đánh giá tác động dài hạn lên kỹ năng tư duy phản biện và gỡ lỗi của sinh viên.
\end{itemize}
